\scanpage{29}
\mspage{10}
\margin{\uline{NB} Ce sous-groupoïde $\mathrm{III}_{J,i}$ (de groupe fondamental $\widehat{\mathbb Z}$) engendré par les sommets $Q_j$ $(j\in J\setminus\{i\})$ et les deux $\unsure{\widetilde a_i^\omega}$ \unsure{orientés} qui les joignent, \uuline{n'est \deleted{\unsure{donc}} pas} stable \uline{par $u$} !}
\unsure{en|si} d'autres termes faisons opérer $u$ sur $\pi(Q_{\omega'i})$, identifié à $\widehat\Pi_{0,3}$ via $\psi_{\omega'i}^{\omega'}$, $\widetilde\lambda_i^\omega$ identifié à $l_0$, on trouve que $u(l_0)$ est conjugué de $l_0^p$, \deleted{avec} la conjugaison étant effectuée par $\eta^{-1}$ (par la formule $(30)$).

Le passage de $u$ à $\tau(u)$ s'exprime par
\deleted{J'en prends maintenant les points $P_\omega^{\unsure{\omega}}$ $(\omega\in\uline{\omega})$}
le passage de $\widetilde\alpha$ à
\[
(31)\qquad\widetilde\alpha'=\varepsilon_0\tau_\infty(\widetilde\alpha^{-1})
=\underbrace{\rho^{-1}\widetilde\sigma_\infty}_{\widetilde\sigma_0}(\widetilde\alpha^{-1})
\]
\margin{où $\widetilde\sigma_\infty=\sigma_\infty\tau_\infty=\tau_\infty\sigma_\infty$ est dans l'image de l'hom canonique
\diagram{\begin{tikzcd}[ampersand replacement=\&]\widetilde G_3\arrow[r]\&\widetilde\Pi_{0,3}\end{tikzcd}}
$\simeq$\quad\deleted{\unsure{du}\ $\widetilde G_3$}
stabilisateur de $P_+$ dans $G_3\times\{1,\tau\}$}
\[
\widetilde\sigma_i(l_i)=l_i^{-1}
\]
\[
\widetilde\sigma_i(l_j)=l_k^{-1}\quad\text{si }\{0,1,\infty\}=\{i,j,k\}
\]
\uline{NB} le groupe $\widetilde G_3$ opère sur $\Pi_{0,3}$ et contient \deleted{$\unsure{\xi}$} \deleted{\unsure{dont il}} $\{1,\rho^{-1},\rho^{-2}\}$ comme ss-groupe \inserted{$\widetilde G_3^+$} d'indice 2, \deleted{\unsure{dont $\widetilde G_3$ conjugue $\rho$}} $\widetilde G_3/\widetilde G_3^+$ opère sur $\widetilde G_3^+$ en échangeant $\rho$ et $\rho^{-1}$, il s'ensuit que $\widetilde G_3$ opère sur $V$ ens des solutions $x$ de l'équation
\[
x\rho^{-1}(x)\rho^{-2}(x)=1
\]
par
\[
(32)\qquad\Theta_g x=g.x^{\unsure{\operatorname{sg}(g)}}
\]
et on a ici
\[
(33)\qquad\widetilde\alpha'=\Theta_{\widetilde\sigma_0}(\widetilde\alpha)
\]
\margin{28}
